\textbf{Desarrollo de las magnitudes intermedias.} Con la aproximación indicada $L\simeq W$, obtener primero CL y resistencia de la polar. La potencia útil disponible es rendimiento por potencia al eje; solo su exceso respecto a $DV$ puede transformarse en altura.\begin{enumerate}

\item Formar fuerza de referencia.

\[\fitmath{\mathrm{qS}=0.5\cdot 1.02\cdot {59}^{2}\cdot 21.4\approx 37991.63\,\mathrm{N}}\]

\item Imponer sustentación aproximada igual a peso.

\[\fitmath{C_L=\frac{2100\cdot 9.80665}{\mathrm{qS}}}\]

\[\fitmath{C_L=\frac{2100\cdot 9.80665}{(37991.634)}\approx 0.5420658\,}\]

\item Calcular resistencia.

\[\fitmath{D=\mathrm{qS}\cdot \left(0.028+0.049\cdot {C_L}^{2}\right)}\]

\[\fitmath{D=(37991.634)\cdot \left(0.028+0.049\cdot {(0.5420657874)}^{2}\right)\approx 1610.767\,\mathrm{N}}\]

\item Potencia aerodinámica necesaria.

\[\fitmath{\mathrm{Pn}=D\cdot 59}\]

\[\fitmath{\mathrm{Pn}=(1610.766661)\cdot 59\approx 95035.23\,\mathrm{W}}\]

\item Convertir potencia al eje en potencia útil.

\[\fitmath{\mathrm{Pd}=0.78\cdot 230000\approx 179400\,\mathrm{W}}\]

\item El exceso de potencia se divide por peso, no por masa.

\[\fitmath{\mathrm{ROC}=\frac{\mathrm{Pd}-\mathrm{Pn}}{2100\cdot 9.80665}}\]

\[\fitmath{\mathrm{ROC}=\frac{(179400)-(95035.23299)}{2100\cdot 9.80665}\approx 4.096577\,\mathrm{m/s}}\]

\end{enumerate}
